% !TEX program = xelatex
% =============================================================================
%  全国大学生数学建模竞赛（CUMCM）模板 —— 电子版
%
%  排版引擎：XeLaTeX（编译两遍以解决交叉引用）
%  编译：在 paper/ 目录内  xelatex -interaction=nonstopmode main.tex  （跑两遍）
%
%  本文件只管赛事差异（文档类、封面/摘要页约定、章节骨架、附录编号口径）。
%    所有通用排版规范在 _base/ 里，由下面的 \input 引入：
%      _base/preamble.tex   页面/字体/行距/标题/caption/表格/浮动体/代码
%      _base/macros.tex     \papertitle \abstractcn \referencescn \appendixAcn
%                           \threelinetable \paperfigure \paperfigurepair…
%      _base/math-style.tex 数学字体与环境
%    9Paper-writing 复制模板时必须把 _base/ 一并复制到 paper/_base/。
%
%  格式规则是机械保证的（写在 _base 里），不是靠写作者自觉。
%    format 阶段独占 _base/ 与本文件导言区的所有权；write 只写 sections/ 和摘要。
%    「正文所有东西的间距、距离上下行、大小都一样，标题统一」——这条
%    由 preamble 保证，不要在每个 section 里手工 \vspace 调间距。
% =============================================================================
\documentclass[12pt, a4paper]{ctexart}

\input{_base/preamble.tex}
\input{_base/macros.tex}

% ---------- 结果登记值（由 `python -m lib.result_contract export` 生成，禁止手工编辑） ----------
% 正文里的数值一律写 \ResultValue{<指标ID>}，不写裸数字 —— 这样改结果时只需重建登记表，
% 不会出现「正文数字与 results/*.json 不一致」。文件不存在时给一条可读的提示而不是报错。
\IfFileExists{result-values.tex}{\input{result-values.tex}}{}
% 未生成时给一条可读提示而不是致命错误。\providecommand 在文件已定义 \ResultValue 时不覆盖。
% 必须用 \detokenize：#1 是指标 ID，而本项目所有真实 ID 都带下划线
%   （q1_surface_moisture_1800s…）。直接排进文本模式会因 `_` 触发
%   `! Missing $ inserted` + `! Extra }, or forgotten $` —— 也就是这个"可读提示"
%   恰好在它唯一的用途（result-values.tex 尚未生成的早期写作）下变成硬编译错误。
\providecommand{\ResultValue}[1]{\textbf{[未生成 result-values.tex：\detokenize{#1}]}}

% ---------- 断行与孤行控制 ----------
\clubpenalty=10000
\widowpenalty=10000
\pretocmd{\section}{\Needspace{8\baselineskip}}{}{}
% 附录里定义的图号（图 A1、A2…）若正文要提前引用，把 \newlabel 桩放在 labels_appendix.tex，
% 由本行引入。没有该文件时静默跳过（多数题目不需要）。
% 两种命名都支持：下划线写法 `labels_appendix.tex` 与连字符写法 `labels-appendix.tex`。
% 桩文件名写错会被 IfFileExists 静默跳过，
% 后果是正文里提前引用的附录图号全变成 ??，而且只留一条 LaTeX 警告。
\IfFileExists{labels_appendix.tex}{\input{labels_appendix.tex}}{%
  \IfFileExists{labels-appendix.tex}{\input{labels-appendix.tex}}{}}

\begin{document}
\hypersetup{pdftitle={},pdfauthor={}}

% =============================================================================
%  标题 + 摘要（电子版第一页；摘要限 1 页；关键词单独一行）
% =============================================================================
\papertitle{<论文标题>}

\abstractcn{%
% 结构：总起一句（本文做了什么）→ 逐问一段（方法 + 精确数值）→ 收尾一句（评价/推广）。
% 每问一段必须以「针对问题N，」开头。数值必须来自 RESULTS_REPORT.md，不得编造。
% 篇幅：约 801 字用正文行距正好排满首页，超过就会溢到第 2 页（那是硬错误）。
<摘要正文：总起句。>

\textbf{针对问题一，}<方法一句话 + 关键结果数值。>

\textbf{针对问题二，}<方法一句话 + 关键结果数值。>

\textbf{针对问题三，}<方法一句话 + 关键结果数值。>
%
}{<关键词1>；<关键词2>；<关键词3>；<关键词4>}

% =============================================================================
%  正文
% =============================================================================
\section{问题重述}

\subsection{问题背景}
% 用自己的话概述题面背景（不要把题目原文整段抄下来）。篇幅控制在半页以内。

\subsection{问题提出}
% 逐条列出题面明确编号的顶层问题。只列编号问题，不要把背景描述、数据说明、提交要求
% 误当成子问题。子问题数按题面实际，不硬凑三个。

\textbf{问题 1：}

\textbf{问题 2：}

\textbf{问题 3：}

% 图 1：分析流程图（`fig_roadmap`）排在这里 —— 问题重述之后、问题分析之前。
%   参考件 `1_restatement.tex` 第 40 行就是这一句，
%   一字不差：`\paperfigure[0.85\textwidth]{fig_roadmap}{分析流程图}{fig_roadmap}`。
%   全文唯一的大图、唯一带配色的图；其余流程图/示意图一律半栏以内、黑白。
%   不得挪进 2_analysis.tex；**别缩，0.85 是下限** —— 954×1297 的竖版大图缩到
%   七成栏宽，节点字直接糊掉。

% 下面这句**必须取消注释**。模板里留成注释，是因为模板目录下没有 figures/*.pdf，
%    放开会让模板本身编译不过（`! Unable to load picture`）。
%    `python lib/web/check_skeleton.py` 会核它有没有放开 —— 放开是硬要求，别忘。
% \paperfigure[0.85\textwidth]{fig_roadmap}{分析流程图}{fig_roadmap}

\section{问题分析}
% 篇幅：**全节 ≤ 一页半，每个问题的分析 ≤ 3 行，每条创新点 3~5 行**。先一句总结性话语，再逐问分析，最后单列创新点。
%   标尺是参考件 `基于变物性耦合传热传质与动域模型的.pdf`：它的 2.1~2.4 各占 2~3 行、
%     创新点每条 4~5 行、全节 1.23 页。本模板排版下 1 行 ≈ 38 字、1 页 ≈ 989 字、一页半 ≈ 1480 字。
%   硬指标：每个「问题N的分析」≤ 3 行（≈114 字）；每条「创新点X：」**3~5 行**（≈114~190 字，
%     逐条量，三条各 5 行是正确写法；下限也是硬的 —— 只写一句时所有上限都满足，
%     于是"重写过"与"把上一版原样拿来用"分不出来）；全节 ≤ 1480 字（一页半），1250 字以上提醒。
%   **「一个数字都不写」只约束「问题N的分析」**：特征时间、下降百分比、倍数、量级、
%     参数明细全是结果，归第五节求解与第六节检验 —— 写在这里既超页、又和后面重复。
%     这里只回答三件事：这一问要什么、难在哪、走什么路线。
%   创新点可以引一个关键数字当证据（参考件每条都带一个）—— 它写的是
%     「做了什么 + 带来什么实质改进」，比问题分析长是对的。
%   `python lib/web/check_skeleton.py` 会逐条核（全节超一页半、或某节超行数上限都判 FAIL，
%     并列出分节明细告诉你该砍哪一段）。

% 下面「问题一/二/三的分析」是举例，不是固定三节：
%    题面有几个子问题，就写几节分析，一节不缺、一节不多。
%     例：题面四问就必须写到「问题四的分析」，漏掉任何一问都算漏节。
%     子问题数在 `5_problemN.tex` 的文件数上——建/删文件时同步建/删这里的分析节。
%     `python lib/web/check_skeleton.py` 会核这一条。

\subsection{问题一的分析}

\subsection{问题二的分析}

\subsection{问题三的分析}

\subsection{本文主要创新点}
% 逐条写，每条 **4~5 行**（≈190 字）：写清「做了什么 + 带来什么实质简化/改进」，
%   并引一个关键数字当证据（参考件每条都带一个）。不要写成方法罗列，也不要只写一句
%   —— 只写一句会短到 2 行，撑不起这一节。逐条独立成段，段间空行。
%   参考件是 3 条（每条 ≈178 字）；4 条时全节约 1215 字 ≈ 1.23 页，仍在一页半以内
%     —— **4~5 条也放得下**，不必为省页数把每条压成一句。

\textbf{创新点一：}

\textbf{创新点二：}

\textbf{创新点三：}

\section{模型假设}
% 篇幅上限：一页；条数 **3--5 条（最多 5 条）**，每条 1.5--2 行（约 55--80 字）。
% 假设节只收「建模约定」。题面已给定的条件、物理/数学必然（如「高度>0 才能起爆」）
%    不放进本节 —— 那是题面重述，放了要判返修。判据口径移到正文判据定义处。
%    每条假设末尾用一句话说明「放宽它会影响哪一问」。
%
% 排面（照参考件 `框架模板.pdf` 第 3 页「三、模型假设」的样子）：
%   ① 标题下先写一句引入语「为了适当地对模型进行合理简化，本文给出如下假设：」，
%      这句顶格（用 \noindent，不留两格缩进）。
%   ② 条目编号用 「假设 N：」，**不要项目符号**（不要 `•`）——
%      `[label=假设 \arabic*：]`；条目标签缩进两格、折行顶格（与正文段落同款）。
%   ③ 条目之间不加空距：与正文段落之间的距离一致（全局 `\setlist` 已把 itemsep/parsep 清零）。
%   参考件量值：引入句 x0=94.9、条目 x0=102.7、折行 x0=88.9、条目间距 19.8pt（= 正文行距）。

\noindent 为了适当地对模型进行合理简化，本文给出如下假设：

\begin{enumerate}[label=\textbf{假设\chinese{enumi}：}]
  \item <建模约定一>。放宽影响：<哪一问、怎么变>。
  \item <建模约定二>。放宽影响：<…>。
  \item <建模约定三>。放宽影响：<…>。
\end{enumerate}

\FloatBarrier
\section{符号说明}
% 篇幅上限：一页。
% 闭环要求（15Verification 会当硬错误抓）：正文用到的每个符号都在本表里；本表每条都在正文出现过。
%
% 本表三列：符号 / 含义 / 单位，单位单独成列，不塞进「含义」的括号里。
%    （参考件 `基于变物性耦合传热传质与动域模型的.pdf` 即此形制。）
%    写成 `$\rho$ & 湿基密度（kg/m$^3$）` 是错的，必须写成两格：
%    `$\rho$ & 湿基密度 & kg/m$^3$`。同一行有多个量时，单位列用分号并列，与含义列逐项对应：
%    `$c_p,\ k$ & 比热容、热传导系数 & J/(kg$\cdot$K)、W/(m$\cdot$K)`；
%    无量纲写「无量纲」，不写「-」。
%
% **本节不写开头总述段**（不要开头总结）：
%   如「除特别说明外，温度用摄氏度…全文所用符号的含义与单位汇总于表 1。」这种段——
%   那是总述/元话语，不是符号说明本身；读者要的是表。单位口径若确实需要交代，
%   写进表注或正文该量出现处，不在这里铺一段。
%   直接 \section{符号说明} 之后出表（下面那段注释讲的是表体本身怎么排）。

% 本表不作浮动体：它是本节唯一的实体内容，高约 20cm，
% 用 table 浮动环境时会因当页放不下而漂到下一页——会漂到「五、模型的建立与求解」
% 标题之后，看上去像符号表跑进了第 5 节。
% 也不用 center+tabular：那样整表不可分割，而本页顶部常被图占去约 207pt，
% 剩余空间装不下整表，TeX 只能把整张表推到下一页，本节标题与正文之后留下约半页空白。
% 改用 longtable：表体可在行间断页，页面填满且始终留在本节之内。
\setlength{\LTleft}{\fill}
\setlength{\LTright}{\fill}
% 必须显式补一段竖向空白：preamble 为统一行距设了 \lineskip=0pt、\lineskiplimit=-20pt，
% 而 longtable 是直接进竖向列表的盒子，TeX 会在它与上一段最后一行之间插入行间胶，
% 该胶在上述设置下被压成 0pt，表头会直接压在「…汇总于表 1。」那一行上。
% 显式 \vspace 之后上一项变成胶，TeX 不再插行间胶，间距恢复正常。
\par\vspace{0.25\baselineskip}
% 表标题放在 longtable 之外，作为一段居中小字（字号与 \caption 的 capstyle 一致）。
% 放进表内当一行时，该行行高由 \@arstrutbox 决定、又受 \lineskip=0pt 影响，
% 只比下一行高 5.5pt，标题被压在 \midrule 上、与表头重叠。移出表外即可完全避开。
% \nopagebreak 保证标题不会与表体分处两页。
{\centering\fontsize{10.5pt}{12.6pt}\selectfont\bfseries 表 1\quad 符号说明\par}
\nopagebreak
\vspace{0.5em}
{\small
  % 行距补丁：preamble 全局设了 \lineskip=0pt、\lineskiplimit=-20pt（为统一正文行距）。
  % 这套设置在普通 tabular 里没问题，但 longtable 会因此把每一行的行间胶压成 0，
  % 表头行只剩 5.6pt 高、\toprule 横穿「符号 含义 单位」。表内改回正规矩即可；
  % \lineskip 取 \baselineskip，即使回退到它也等于正常行距，不会压行。
  \setlength{\lineskiplimit}{0pt}%
  \setlength{\lineskip}{\baselineskip}%
  % 表头行由 5.6pt 恢复到正常高度后整表高了约 20pt，本表行距单独收到 1.0 以抵消。
  \renewcommand{\arraystretch}{1.0}%
  % 三列对齐：符号靠左、含义居中、单位靠右
  \begin{longtable}{@{}L{0.19\textwidth}C{0.57\textwidth}R{0.19\textwidth}@{}}
    % 表头行不能放进 \endfirsthead：那样 \toprule 会被画到「符号 含义 单位」中间，
    % 且表头下的 \midrule 直接丢失——这是 longtable 首块的手工规则放置问题。
    % 把表头行改成正文首行即可两者兼得：横线位置正确，且续页不重复表头。
    \endfirsthead
    \endhead
    \bottomrule
    \endlastfoot
    \toprule
    符号 & 含义 & 单位 \\
    \midrule
    % ↓↓ 以下两行只是版式示例（照抄会判「符号表与正文对不上」）：全部替换为本题符号。
    %    排列次序照正文出现次序或按类别分组（几何/物性/环境/数值），别按字母表排。
    %    第一行示「一符一义」；第二行示「一格多个量」——含义与单位用顿号/分号逐项对齐。
    %    占位符用 ASCII（`$<S_1>$`）是有意的：汉字进数学模式渲染不出来，留空会看不清版式。
    $<S_1>$ & <含义：只写「是什么」，不写单位> & <单位> \\
    $c_p,\ k$ & 比热容、热传导系数 & J/(kg$\cdot$K)、W/(m$\cdot$K) \\
  \end{longtable}}

% 收尾（可选但常拿分）：把同名不同义的量点清楚，凡有歧义就写这一段。
% 例：「附件 1 的 $C_\infty$ 是空气含湿量，与药材干基含水率 $C$ 基准不同（见模型假设）」
%     「$\bar C_\infty$（平台环境均值）、$C_{lo}$（环境下确界）与 $\bar C$（药材体积平均）三者互不相同」。


\section{模型的建立与求解}
% 每问统一小结构（评审偏好）：具体分析 → 模型准备 → 模型建立 → 模型求解 →（选）结果分析。
% 用连贯学术段落叙述，少碎标题、少列表。问题一不要写得特别长：序号越大按理内容越多。
\subsection{问题一的模型的建立和求解}

\subsubsection{具体分析}

\subsubsection{模型准备}

\subsubsection{模型建立}
% 决策变量 / 目标函数 / 约束条件（建模三要素）在这里显式写出。
% 数值一律用 \ResultValue{<指标ID>}，不写裸数字。

\subsubsection{模型求解}

\subsection{问题二的模型的建立和求解}

\subsubsection{具体分析}

\subsubsection{模型准备}

\subsubsection{模型建立}

\subsubsection{模型求解}

\subsection{问题三的模型的建立和求解}

\subsubsection{具体分析}

\subsubsection{模型准备}

\subsubsection{模型建立}

\subsubsection{模型求解}

\subsubsection{结果分析}
% 问题三通常最难、内容最多，结果分析往往不可省。

% 题目不止三问时按同样的命名往后加 5_problem4…；子问题数由题面决定，不硬凑。

\section{模型检验与分析}
% 篇幅上限：两页。内容来自 reports/ROBUSTNESS_REPORT.md 的「给 9Paper-writing 的可直接引用素材」。
% 灵敏度章节**至少要有 1 张可视化图**，整节纯数字表会被判 WARN。

\subsection{误差分析}

\subsection{灵敏度分析}

\section{模型的评价与推广}
% 篇幅上限：两页。
% 「优点」不得与「缺点」「灵敏度」节逐字重复 ——
%    写之前先对照 6_check.tex，避免同一段话在两处出现。

\subsection{模型的优点}

\subsection{模型的缺点}

\subsection{模型的改进与推广}


% =============================================================================
%  八、AI 工具使用说明；九、参考文献（合计不超过 1 页）
%  这里不加 \clearpage —— 强令另起一页会把参考文献与附录各顺延一页。
% =============================================================================
\section{AI 工具使用说明}
本参赛队在竞赛过程中使用了 AI 工具，主要用于文献搜集、语言美化、图表优化、代码测试，
详细使用情况见支撑材料。

\referencescn

% =============================================================================
%  附录（正文的最后一页）
%
%  只放一张「支撑材料文件清单」表：逐项登记提交出去的文件名 + 它做什么用。
%  表体由 14Layout-and-format 阶段据 reports/SUBMISSION_MANIFEST.json 生成 —— 那份清单同时是
%  打包依据，两边必须逐项一致（`python -m lib.delivery check` 会对账）。
%
%  详细的补充推导 / 补充图表不在这里，它们单独编译成 `附录A.pdf`
%     （工程在 `paper_appendix/`，模板见 `appendix/`）。两个 PDF 的公式编号要接续，
%     起始编号写在本行之前的 `\setcounter{equation}{N}`（N = 正文公式总数）。
%
%  这里没有「附录 B 完整源程序」：源码不再排进论文 ——
%     代码以散装 `.py` 提交，在本表里逐项登记即可。
% =============================================================================
\appendixAcn[支撑材料清单]{A1_materials}

\end{document}
